%% Font size %%
\documentclass[11pt]{article}

%% Load the custom package
\usepackage{Mathdoc}

%% Numéro de séquence %% Titre de la séquence %%
\renewcommand{\centerhead}{Fractions - Éval 1 : Coloriages}

%% Spacing commands %%
\renewcommand{\baselinestretch}{1} \setlength{\parindent}{0pt}

%% Carré divisé en #1 parts égales (bandes verticales), #2 parts coloriées, côté #3 cm
\newcommand{\squarefrac}[3]{%
\begin{tikzpicture}[baseline=0pt]
\ifnum#2>0
\foreach \i in {0,...,\numexpr#2-1\relax}{
\fill[gray!55] ({\i*#3/#1},0) rectangle ({(\i+1)*#3/#1},#3);
}
\fi
\foreach \i in {1,...,\numexpr#1-1\relax}{
\draw (\i*#3/#1,0) -- (\i*#3/#1,#3);
}
\draw[thick] (0,0) rectangle (#3,#3);
\end{tikzpicture}%
}

%% Camembert divisé en #1 secteurs égaux, #2 secteurs coloriés, rayon #3 cm
\newcommand{\piefrac}[3]{%
\begin{tikzpicture}[baseline=0pt]
\ifnum#2>0
\foreach \i in {0,...,\numexpr#2-1\relax}{
\fill[gray!55] (0,0) -- ({90-\i*360/#1}:#3) arc ({90-\i*360/#1}:{90-(\i+1)*360/#1}:#3) -- cycle;
}
\fi
\foreach \i in {0,...,\numexpr#1-1\relax}{
\draw (0,0) -- ({90-\i*360/#1}:#3);
}
\draw[thick] (0,0) circle (#3);
\end{tikzpicture}%
}

%% Zone de réponse bien visible, sous forme de fraction (barre horizontale, comme \dfrac)
\newcommand{\reponsefrac}{%
\textbf{Fraction :} \ $\dfrac{\ldots\ldots}{\ldots\ldots}$%
}

%% Une question, dans une boîte de largeur fixe (le \\ interne reste bien local à la boîte)
\newcommand{\itembox}[1]{\begin{minipage}[t]{8.3cm}\centering #1\end{minipage}}

%% Une ligne de la grille : une question à gauche, une à droite
\newcommand{\ligne}[2]{\noindent\itembox{#1}\hfill\itembox{#2}\par\vspace{.7cm}}

\begin{document}

\phantom{0}
\vspace{-.5cm}

\entetedevoirs{20}

\begin{center}
\duree{35 minutes}
\total{20 points}
\soin
\end{center}

\begin{exercicedevoir}[5][Colorier une fraction (fraction < 1)]
Pour chaque figure, colorie la fraction indiquée.

\ligne{(a) Colorie $\dfrac{1}{2}$\\[.2cm]\squarefrac{2}{0}{1.8}}{(d) Colorie $\dfrac{3}{5}$\\[.2cm]\piefrac{5}{0}{1.3}}
\ligne{(b) Colorie $\dfrac{2}{3}$\\[.2cm]\piefrac{3}{0}{1.3}}{(e) Colorie $\dfrac{5}{6}$\\[.2cm]\squarefrac{6}{0}{1.8}}
\ligne{(c) Colorie $\dfrac{3}{4}$\\[.2cm]\squarefrac{4}{0}{1.8}}{(f) Colorie $\dfrac{1}{4}$\\[.2cm]\piefrac{4}{0}{1.3}}
\end{exercicedevoir}

\begin{exercicedevoir}[5][Retrouver une fraction (fraction < 1)]
Pour chaque figure déjà coloriée, écris la fraction correspondante dans les cases.

\ligne{(a) \piefrac{2}{1}{1.3}\\[.2cm]\reponsefrac}{(d) \squarefrac{5}{4}{1.8}\\[.2cm]\reponsefrac}
\ligne{(b) \squarefrac{3}{2}{1.8}\\[.2cm]\reponsefrac}{(e) \piefrac{6}{5}{1.3}\\[.2cm]\reponsefrac}
\ligne{(c) \piefrac{4}{1}{1.3}\\[.2cm]\reponsefrac}{(f) \squarefrac{4}{3}{1.8}\\[.2cm]\reponsefrac}
\end{exercicedevoir}

\begin{exercicedevoir}[5][Colorier une fraction (fraction > 1)]
Pour chaque fraction, colorie le nombre de parts correspondant sur les figures proposées (une figure entière correspond à 1).

\ligne{(a) Colorie $\dfrac{3}{2}$\\[.2cm]\piefrac{2}{0}{1.3}\ \piefrac{2}{0}{1.3}}{(d) Colorie $\dfrac{8}{5}$\\[.2cm]\squarefrac{5}{0}{1.6}\ \squarefrac{5}{0}{1.6}}
\ligne{(b) Colorie $\dfrac{5}{3}$\\[.2cm]\squarefrac{3}{0}{1.6}\ \squarefrac{3}{0}{1.6}}{(e) Colorie $\dfrac{9}{6}$\\[.2cm]\piefrac{6}{0}{1.3}\ \piefrac{6}{0}{1.3}}
\ligne{(c) Colorie $\dfrac{7}{4}$\\[.2cm]\piefrac{4}{0}{1.3}\ \piefrac{4}{0}{1.3}}{(f) Colorie $\dfrac{5}{4}$\\[.2cm]\squarefrac{4}{0}{1.6}\ \squarefrac{4}{0}{1.6}}
\end{exercicedevoir}

\begin{exercicedevoir}[5][Retrouver une fraction (fraction > 1)]
Pour chaque ensemble de figures déjà coloriées, écris la fraction correspondante dans les cases (une figure entière correspond à 1).

\ligne{(a) \piefrac{2}{2}{1.3}\ \piefrac{2}{1}{1.3}\\[.2cm]\reponsefrac}{(d) \squarefrac{5}{5}{1.6}\ \squarefrac{5}{3}{1.6}\\[.2cm]\reponsefrac}
\ligne{(b) \squarefrac{3}{3}{1.6}\ \squarefrac{3}{2}{1.6}\\[.2cm]\reponsefrac}{(e) \piefrac{6}{6}{1.3}\ \piefrac{6}{3}{1.3}\\[.2cm]\reponsefrac}
\ligne{(c) \piefrac{4}{4}{1.3}\ \piefrac{4}{3}{1.3}\\[.2cm]\reponsefrac}{(f) \squarefrac{4}{4}{1.6}\ \squarefrac{4}{1}{1.6}\\[.2cm]\reponsefrac}
\end{exercicedevoir}

\end{document}

%%% Local Variables:
%%% mode: LaTeX
%%% TeX-master: t
%%% End:
